\subsection{Removing exceptional centers}
\label{subsec:repair}

The pointwise probability bound will control the measure of the centers
at which some normalized scale is missed.  We first enlarge the random
set to an open set, so that those exceptional centers form a closed
set.  A small open neighborhood of that closed set will then provide
hits from the dyadic tail.
This follows the open-cover completion of Kolountzakis
\cite[Sections~2.2 and~3.2]{Kolountzakis1997};
we give the closedness, measure and tail arguments explicitly.

\begin{proof}[Completion of the proof of Lemma~\ref{lem:periodic}]
The probability space \(\Omega\) of all table entries is finite.
For each outcome \(\omega\), the set \(B_\omega\) is a union of cells
from one finite dyadic grid per period: its finest key determines
every table lookup.  Regard a periodic set as a subset of the circle
\(\mathbb T=\mathbb R/\mathbb Z\), with normalized Lebesgue measure
equal to the density of its periodic lift.  Taking the closures of the selected
cells adds only finitely many endpoints.  Enlarging these closed arcs
by sufficiently small open circular neighborhoods therefore gives an
open periodic set \(B_\omega^+\supseteq B_\omega\) with
\begin{equation}
 \rho(B_\omega^+)\le\rho(B_\omega)+p.
 \label{eq:open-enlargement}
\end{equation}
For example, if the grid has \(Q\) cells, enlargement by circular
distance \(p/(4Q)\) adds at most \(p/2\) to the measure of their union.
For the empty set use the empty enlargement.  Since there are only
finitely many outcomes, these choices introduce no measurable-selection
issue.  Equation~\eqref{eq:mean-density} gives
\begin{equation}
 \mathbb E\rho(B_\omega^+)\le2p.
 \label{eq:open-mean}
\end{equation}

Define the exceptional-center set
\begin{equation}
 R_\omega=
 \left\{x\in\mathbb R:
   \exists t\in[1,2]\ \forall n\in\mathcal N,\
          x+t2^{-n}\notin B_\omega^+\right\}.
 \label{eq:exceptional-centers}
\end{equation}
This set is periodic and closed.  To prove closedness without any
endpoint convention on a fundamental interval, let
\(\pi:\mathbb R\to\mathbb T\) be the quotient map and let
\(B_{\omega,\mathbb T}^+\subset\mathbb T\) denote the open image of
\(B_\omega^+\).  Each map
\[
 (\pi x,t)\longmapsto\pi(x+t2^{-n})
 \quad\hbox{from }\mathbb T\times[1,2]\hbox{ to }\mathbb T
\]
is continuous.  The conditions that all these images avoid
\(B_{\omega,\mathbb T}^+\), for \(n\in\mathcal N\), consequently define
a closed subset of the compact space \(\mathbb T\times[1,2]\).
Its projection to \(\mathbb T\) is compact and hence closed.
Its inverse image under \(\pi\) is exactly \(R_\omega\).

In particular, each indicator \({\bf1}_{R_\omega}(x)\) is Borel
measurable in \(x\).  Since \(\Omega\) is finite, it is also jointly
measurable in \((\omega,x)\).  For every stable center \(x\in G\),
the inclusion \(B_\omega\subseteq B_\omega^+\) and
Proposition~\ref{prop:raw-failure} give
\[
 \mathbb P(x\in R_\omega)\le2p.
\]
At the other centers the probability is at most \(1\).
Interchanging a finite sum and an integral, and using
Lemma~\ref{lem:stable}, we obtain
\begin{align}
 \mathbb E\rho(R_\omega)
 &=\int_0^1\mathbb P(x\in R_\omega)\,dx \notag\\
 &\le 2p\,m(G\cap[0,1])+m(G^c\cap[0,1])
 \le2p+\rho(G^c)\le3p.
 \label{eq:exceptional-mean}
\end{align}
Thus \eqref{eq:open-mean} and \eqref{eq:exceptional-mean} imply
\[
 \mathbb E\bigl(\rho(B_\omega^+)+\rho(R_\omega)\bigr)\le5p.
\]
Choose one outcome whose summed density is at most this bound, and
write its sets as \(B^+\) and \(R\).  We have
\begin{equation}
 \rho(B^+)+\rho(R)\le5p.
 \label{eq:chosen-outcome}
\end{equation}

It remains to cover every center in \(R\).  Its image
\(R_{\mathbb T}=\pi(R)\) is closed in the circle.  If it is empty,
put \(V=\varnothing\).  Otherwise, for \(\varepsilon>0\), its open
metric neighborhoods
\[
 V_\varepsilon^{\mathbb T}
 =\{z\in\mathbb T:
           \operatorname{dist}_{\mathbb T}(z,R_{\mathbb T})
                          <\varepsilon\}
\]
decrease to \(R_{\mathbb T}\) as \(\varepsilon\downarrow0\).
Continuity of finite measure from above lets us choose
\(\varepsilon>0\) so that the open periodic lift
\(V=\pi^{-1}(V_\varepsilon^{\mathbb T})\) satisfies
\[
 R\subseteq V,\qquad \rho(V)\le\rho(R)+p.
\]
Set \(H=B^+\cup V\).  It is open and periodic, and
\eqref{eq:chosen-outcome} gives
\[
 \rho(H)\le\rho(B^+)+\rho(V)
          \le\rho(B^+)+\rho(R)+p\le6p.
\]

If \(x\notin R\), negating \eqref{eq:exceptional-centers} shows that
for every \(t\in[1,2]\) some \(n\in\mathcal N\) satisfies
\(x+t2^{-n}\in B^+\subseteq H\).
If \(x\in R\), choose an integer \(n_0\ge1\) with
\(2^{1-n_0}<\varepsilon\).  For every \(t\in[1,2]\) and every
\(n\ge n_0\), the circle distance satisfies
\[
 \operatorname{dist}_{\mathbb T}
       \bigl(\pi(x+t2^{-n}),R_{\mathbb T}\bigr)
 \le \operatorname{dist}_{\mathbb T}
       \bigl(\pi(x+t2^{-n}),\pi x\bigr)
 \le t2^{-n}<\varepsilon.
\]
Hence \(x+t2^{-n}\in V\subseteq H\).
These indices are taken from the full dyadic tail.  The two cases
establish the required hit for every real \(x\) and every
\(t\in[1,2]\), completing the lemma.
\end{proof}
